\documentclass{article}

\usepackage[utf8]{inputenc}
\usepackage{helvet}
\usepackage{geometry}
\usepackage{xcolor}
\usepackage{eso-pic}
\usepackage{fancyhdr}
\usepackage{parskip}
\usepackage{microtype}
\usepackage[english,ukrainian]{babel}
\usepackage{url}
\usepackage{amsmath}
\usepackage{amssymb}
\usepackage{amsthm}
\usepackage{longtable}
\usepackage{booktabs}
\usepackage{hyperref}
\usepackage{titlesec}

\definecolor{nato}{HTML}{293757}
\definecolor{footergray}{gray}{0.65}

\geometry{a4paper,left=3cm,right=3cm,top=3cm,bottom=3cm}
\renewcommand{\familydefault}{\sfdefault}
\setcounter{secnumdepth}{3}

\DeclareFontFamily{T2A}{phv}{}
\DeclareFontShape{T2A}{phv}{m}{n}{<->ssub * cmss/m/n}{}
\DeclareFontShape{T2A}{phv}{bx}{n}{<->ssub * cmss/bx/n}{}
\DeclareFontShape{T2A}{phv}{m}{it}{<->ssub * cmss/m/it}{}
\DeclareFontShape{T2A}{phv}{bx}{it}{<->ssub * cmss/bx/it}{}

\titleformat{\section} {\color{nato}\Huge\bfseries}{\thesection.}{1em}{}
\titlespacing*{\section}{0pt}{30pt}{12pt}
\titleformat{\subsection} {\color{nato}\LARGE\bfseries}{\thesubsection.}{1em}{}
\titlespacing*{\subsection}{0pt}{20pt}{8pt}
\titleformat{\subsubsection} {\color{nato}\Large\bfseries}{\thesubsubsection.}{1em}{}
\titlespacing*{\subsubsection}{0pt}{10pt}{6pt}

\fancyhf{}
\fancyhead[R]{\small\color{footergray} 27 червня 2026}
\fancyfoot[L]{\small\color{footergray} ERP/1. Модуль Конференцзв'язок}
\fancyfoot[R]{\small\color{footergray}\thepage}

\newcommand\HUGE[1]{\fontsize{40}{40}\selectfont #1}
\renewcommand{\headrulewidth}{0pt}
\renewcommand{\footrulewidth}{0.4pt}
\renewcommand{\footrule}{\color{footergray}\hrule width \textwidth height 0.4pt}

\usepackage{amsmath}
\usepackage{amssymb}
\usepackage{booktabs}
\usepackage{hyperref}
\usepackage{tikz}
\usetikzlibrary{shapes.geometric, arrows, arrows.meta, positioning, fit, backgrounds}

\let\originalfootnotesize\footnotesize
\let\oldverbatim\verbatim
\let\oldendverbatim\endverbatim
\renewenvironment{verbatim}{\originalfootnotesize\oldverbatim}{\oldendverbatim}

\title{Architecture of GStreamer WebRTC MCU Media Compositor for Chat Applications}
\author{Namdak Tonpa, Zen Crypted}
\date{\today}

\begin{document}

\begin{titlepage}
  \AddToShipoutPictureBG*{\AtPageLowerLeft{\color{nato}\rule{\paperwidth}{\paperheight}}}
  \color{white}\thispagestyle{empty}\vspace*{3cm}\centering
  {\Huge  \textbf{SYNRC: OpenWebRTC} \par} \vspace{1.5cm}
  {\HUGE  \textbf{Технічна архітектура} \par} \vspace{1.5cm}
  {\LARGE \textbf{Аудіо-відео композитора медіа-потоків для WebRTC з використанням GStreamer} \par} \vspace{1.5cm}
  {\large На основі Технічних Вимог \par}
  {\large Згідно з Постановою КМУ № 205 від 21.02.2025 \par}
  {\large Формалізованих за стандартом ISO/IEC/IEEE 29148:2018 \par}
  \vfill
  {\large 2026 \copyright\ Криптографічні Телесистеми \par}
\end{titlepage}

\newpage
\maketitle

\begin{abstract}
This document provides a specification of the GStreamer-based Multipoint Control Unit (MCU)
implemented in \texttt{c\_src/gst.c} in 1K LOC. It defines the pipeline architecture, data structures,
dynamic peer lifecycle, inter-process communication protocol, output multiplexing strategies,
and architectural invariants of the production C99 compositor binary \texttt{priv/gst}.
\end{abstract}

\newpage
\tableofcontents

\newpage
\section{Introduction}

The `gst` binary is a self-contained C99 process constituting the media plane
of the RTP video conferencing system. It fulfils three simultaneous responsibilities:

\begin{enumerate}
\item \textbf{Upstream Ingest} --- Receives encrypted WebRTC SRTP/SRTCP streams from each
   participant via independent \texttt{webrtcbin} elements, decoding them into raw
   audio-visual frames using per-peer \texttt{decodebin} instances.
\item \textbf{Mixing and Compositing} --- Composites all decoded video feeds into a single
   1920x1080 spatial grid via GStreamer's \texttt{compositor} and additively mixes
   all audio tracks through \texttt{audiomixer}.
\item \textbf{Downstream Broadcast and Recording} --- Re-encodes and broadcasts the composite
   stream back to every participant as a single WebRTC downstream connection while
   simultaneously writing the session to disk in HLS or fragmented MP4 format.
\end{enumerate}

The binary is spawned as a supervised OS port process by the Erlang
\texttt{rtp\_broker} \texttt{gen\_server}. All signaling is exchanged as newline-delimited
JSON over UNIX standard I/O pipes, making the media plane fully decoupled from
the Erlang control plane.

\section{Mixer Architecture}

\subsection{High-Level Architecture}

The following diagram provides a macro-level conceptual overview of the media pipeline.
It illustrates the four primary lifecycle stages (Ingest, Central Mixer, Broadcast, and Recording)
without getting bogged down in the exact pad connections or intermediate conversion elements.
Use this to understand the general data flow of the MCU.

\begin{verbatim}
flowchart TD
    Peer[WebRTC Peer] <--> WB[webrtcbin]
    subgraph Ingest ["Low-Latency Ingest"]
        WB   --> DB[decodebin]
        DB   --> VC[videoconvert]
        VC   --> VR[videoscale+rate+caps]
        VR   --> JQ[v_jitter queue leaky 500ms]
        JQ   --> COMP[compositor Grid]
        DB   --> AC[audioconvert]
        AC   --> AJQ[a_jitter queue leaky 500ms]
        AJQ  --> AMIX[audiomixer]
    end
    subgraph Mixer ["Central Mixer"]
        COMP --> ENC[x264enc ultrafast zerolatency]
        ENC  --> TEE[tee]
        AMIX --> ATEE[audio tee]
    end
    subgraph Broadcast ["Broadcast"]
        TEE  --> VQ[per-peer v_queue leaky 1s] --> WB
        ATEE --> AQ[per-peer a_queue leaky 1s] --> WB
    end
    subgraph Recording ["Recording"]
        TEE  --> MUX[mp4mux / hlssink2] --> OUT[Output File]
        ATEE --> MUX
    end
    classDef critical fill:#FF6B6B,stroke:#FF0000,color:white
    classDef queue fill:#4ECDC4,stroke:#45B7D1
    classDef mixer fill:#45B7D1,stroke:#2C8C9E
    class JQ,AJQ critical
    class VQ,AQ queue
    class COMP,ENC,AMIX mixer
\end{verbatim}

\subsubsection{Low-Level Detailed Topology}

The second diagram below provides the micro-level exact element topology.
It explicitly maps every dynamic pad connection (e.g., \texttt{sink\_N}, \texttt{sink\_0}, \texttt{sink\_1})
and dual-tee bifurcations exactly as they are constructed in \texttt{gst.c}.
Use this diagram as the canonical reference when reading or modifying the C source code.

\begin{verbatim}
flowchart TD
    %% External Peers
    Peer[WebRTC Peer\nBrowser] <--> WB[webrtcbin]

    %% Ingest Path (Per Peer)
    subgraph Ingest ["Ingest Path (Per Peer)"]
        WB -->|src_%u   sink| DB[decodebin]
        DB -->|src_%u| PAD[pad-added]
        PAD -->|video pad   sink| VC
        PAD -->|audio pad   sink| AC
        subgraph VideoIngest [Video Ingest]
            VC[videoconvert] -->|src   sink| VS[videoscale]
            VS -->|src   sink| VR[videorate]
            VR -->|src   sink| VCAPS[capsfilter<br>30/1 I420]
            VCAPS -->|src   sink| JQ[v_jitter queue<br>leaky 500ms]
            JQ -->|src| COMP_SINK[compositor.sink_%u<br>Grid Position]
        end
        subgraph AudioIngest [Audio Ingest]
            AC[audioconvert] -->|src   sink| AR[audioresample]
            AR -->|src   sink| AJQ[a_jitter queue<br>leaky 500ms]
            AJQ -->|src| AMIX_SINK[audiomixer.sink_%u]
        end
    end

    %% Central Mixing
    subgraph Mixer ["Central Mixer"]
        COMP[compositor<br>name=mix<br>ignore-inactive-pads=true] -->|src   sink| VCONV[videoconvert → I420]
        VCONV -->|src   sink| ENC[x264enc<br>ultrafast + zerolatency]
        ENC -->|src   sink| HTEE[h264_tee]
        AMIX[audiomixer<br>name=amix<br>ignore-inactive-pads=true] -->|src   sink| ACONV[audioconvert + resample]
        ACONV -->|src   sink| RAW_A[raw_atee]
    end

    %% Broadcast Distribution
    subgraph Broadcast ["Broadcast to All Peers"]
        HTEE -->|src_%u   sink| VTEE[vtee]
        VTEE -->|src_%u   sink| VQ[per-peer v_queue<br>leaky 1s]
        VQ -->|src   sink_0| WB
        RAW_A -->|src_%u   sink| ATEE[atee]
        ATEE -->|src_%u   sink| AQ[per-peer a_queue<br>leaky 1s]
        AQ -->|src   sink_1| WB
    end

    %% Recording
    subgraph Recording ["Recording Output"]
        HTEE -->|src_%u   video_0| MUX[mp4mux / hlssink2]
        RAW_A -->|src_%u   audio_0| MUX
        MUX -->|src| OUT[recording.mp4<br>or HLS segments]
    end

    classDef webrtc fill:#90EE90,stroke:#228B22
    classDef mixer fill:#FFB6C1,stroke:#DB7093
    classDef queue fill:#87CEEB,stroke:#4682B4
    classDef encoder fill:#FFD700,stroke:#DAA520

    class WB webrtc
    class COMP,AMIX mixer
    class JQ,AJQ,VQ,AQ queue
    class ENC encoder
\end{verbatim}

\subsection{Data Structures}

\subsubsection{Per-Participant Context Record}

\begin{verbatim}
typedef struct {
    gchar      *peer_id;      // Unique participant identifier string
    GstElement *webrtc;       // webrtcbin element (upstream + downstream)
    GstElement *v_queue;      // Leaky downstream video broadcast queue
    GstElement *a_queue;      // Leaky downstream audio broadcast queue
    GstPad     *comp_pad;     // Requested compositor sink pad (retained for cleanup)
    GstPad     *amix_pad;     // Requested audiomixer sink pad (retained for cleanup)
    GstElement *v_decodebin;  // Dynamically attached video decodebin
    GstElement *a_decodebin;  // Dynamically attached audio decodebin
    GstElement *v_convert;    // videoconvert inserted between decodebin and scaler
    GstElement *v_scale;      // videoscale for predictable output dimensions
    GstElement *v_rate;       // videorate for duplicating missing frames
    GstElement *v_caps;       // capsfilter enforcing I420 and 30fps
    GstElement *v_jitter;     // leaky 500ms queue acting as an asynchronous thread boundary
    GstElement *a_convert;    // audioconvert inserted before audioresample
    GstElement *a_resample;   // audioresample normalizing sample rate
    GstElement *a_jitter;     // leaky 500ms audio jitter queue
    gint        grid_idx;               // Grid slot index [0..15] controlling spatial position
    gboolean    remote_desc_set;        // Flag tracking if SDP answer was applied
    gboolean    bundled;                // Flag indicating bundled RTCP
    GArray     *pending_ice_candidates; // ICE candidates awaiting SDP resolution
} PeerInfo;
\end{verbatim}

\texttt{PeerInfo} aggregates every GStreamer element and requested pad dynamically created
on peer join, enabling deterministic full cleanup on departure. All instances are
stored in a \texttt{GHashTable} keyed by \texttt{peer\_id} with value destructor \texttt{free\_peer\_info}.

\subsubsection{Singleton Pipeline State}

\begin{verbatim}
typedef struct {
    GstElement  *pipeline;
    GstElement  *compositor;     // Video compositing mixer (name="mix")
    GstElement  *audiomixer;     // Audio mixing element (name="amix")
    GstElement  *video_tee;      // Broadcast video tee (name="vtee")
    GstElement  *audio_tee;      // Broadcast audio tee (name="atee")
    GHashTable  *webrtcbins;     // peer_id -> PeerInfo*
    GMainLoop   *loop;
    gboolean     grid_slots[16]; // Spatial slot occupancy bitmap (up to 4x4)
    gint         pad_index;      // Global compositor pad counter
} RecorderState;
\end{verbatim}

\texttt{RecorderState} is a process-global singleton. References to the four permanent
pipeline elements are resolved by name after \texttt{gst\_parse\_launch()} and cached
for O(1) dynamic peer attachment.

\subsection{Static Pipeline Topology}

The persistent portion of the pipeline is constructed once at startup via
`gst\_parse\_launch()`. Three variants are supported, selected by the second CLI argument.

\subsubsection{HLS/H.264 (Deafault)}

\begin{verbatim}
videotestsrc pattern=black is-live=true do-timestamp=true !
timeoverlay valignment=bottom halignment=left font-desc=\"Sans, 48\" !
video/x-raw,width=1920,height=1080,framerate=15/1 ! queue max-size-buffers=1 leaky=downstream !
mix.sink_0 audiotestsrc is-live=true do-timestamp=true volume=0 !
queue max-size-buffers=5 leaky=downstream !
amix.sink_0 compositor name=mix ignore-inactive-pads=true ! videoconvert !
video/x-raw,format=I420,width=1920,height=1080,framerate=15/1 !
queue max-size-time=300000000 max-size-buffers=0 max-size-bytes=0 leaky=downstream !
x264enc bitrate=2000 speed-preset=ultrafast key-int-max=15 tune=zerolatency !
video/x-h264,profile=baseline ! h264parse ! tee name=h264_tee h264_tee. !
queue max-size-time=300000000 max-size-buffers=0 max-size-bytes=0 leaky=downstream !
rtph264pay config-interval=1 pt=96 ! tee name=vtee audiomixer name=amix ignore-inactive-pads=true !
audioconvert ! audioresample ! audio/x-raw,rate=48000,channels=2 ! tee name=raw_atee raw_atee. !
queue max-size-time=300000000 max-size-buffers=0 max-size-bytes=0 leaky=downstream !
opusenc ! rtpopuspay pt=111 ! tee name=atee h264_tee. !
queue max-size-time=300000000 max-size-bytes=0 max-size-buffers=0 leaky=downstream flush-on-eos=true !
hlssink2.video raw_atee. !
queue max-size-time=300000000 max-size-bytes=0 max-size-buffers=0 leaky=downstream flush-on-eos=true !
audioconvert ! audioresample ! audio/x-raw,rate=44100,channels=2 ! avenc_aac ! aacparse !
hlssink2.audio hlssink2 name=hlssink2 async-handling=true location=%s/segment_%" G_GINT64_FORMAT "_%%05d.ts
playlist-location=%s/index.m3u8 target-duration=2 max-files=0 playlist-length=10
\end{verbatim}

\subsubsection{Fragmented MP4}

Identical upstream compositing graph; HLS sink replaced by:

\begin{verbatim}
mp4mux name=mux fragment-duration=1000 streamable=true !
  filesink location=<out_dir>/recording.mp4

h264_tee. ! queue leaky=2 max-size-time=30000000000 ! mux.video_0
raw_atee. ! queue leaky=2 max-size-time=30000000000 !
  audioconvert ! audioresample ! audio/x-raw,rate=44100,channels=2 !
  avenc_aac ! aacparse ! mux.audio_0
\end{verbatim}

\subsubsection{HEVC/HLS}

Dual-encoding via `raw\_vtee`: H.264 for WebRTC participants; H.265 for HLS storage.
Audio follows the same `raw\_atee` bifurcation.

\subsubsection{Invariant Bootstrap Sources}

A live black `videotestsrc` on `mix.sink\_0` (zorder=1) and a silent `audiotestsrc`
on `amix.sink\_0` are permanently active. They prevent scheduler stalls when no peers
are connected and allow dynamic peer sources to preroll instantaneously.

\subsection{Dynamic Peer Lifecycle}

\subsubsection{Peer Join}

Six ordered operations on receipt of \texttt{\{"type":"peer\_joined","peer\_id":"..."\}}:

\begin{enumerate}
\item Grid slot allocation: Scans `state.grid\_slots[0..15]` for the first free slot.
\item \texttt{webrtcbin} creation: Named element with 50 ms latency budget.
\item Leaky broadcast queue creation: `leaky=2`, `max-size-time=1000000000` (1 s).
\item Downstream broadcast linkage: `vtee -> v\_queue -> webrtcbin.sink\_0`; `atee -> a\_queue -> webrtcbin.sink\_1`.
\item State synchronization: `gst\_element\_sync\_state\_with\_parent()` on all added elements.
\item SDP offer generation: Asynchronous `GstPromise` triggers `on\_offer\_created`.
\end{enumerate}

\subsubsection{Incoming RTP Track Routing}

\begin{verbatim}
const gchar *media = gst_structure_get_string(str, "media");
gboolean is_video = (g_strcmp0(media, "video") == 0);
\end{verbatim}

A `decodebin` is linked to the RTP pad; decoded pads trigger `on\_decoded\_pad`.

\subsubsection{Decoded Pad Routing}

\textbf{Video path} --- compositor quadrant placement:

\begin{verbatim}
GstPad *comp_pad = gst_element_request_pad_simple(state.compositor, "sink_%u");
peer->comp_pad = comp_pad;

gint idx = peer->grid_idx;
gint w = WIDTH / 2;   // 960
gint h = HEIGHT / 2;  // 540
gint x = (idx % 2) * w;
gint y = (idx / 2) * h;

g_object_set(comp_pad, "xpos", x, "ypos", y, "width", w, "height", h,
             "zorder", (guint)(idx + 10), "sizing-policy", 1, NULL);
\end{verbatim}

**Audio path**: `audioconvert -> audioresample -> audiomixer.sink\_\%u`.

\subsubsection{Peer Departure — Six-Stage Teardown}

\begin{verbatim}
| Stage | Action                                                                                       |
|-------|----------------------------------------------------------------------------------------------|
| 1     | Unlink/release vtee src pad; NULL v_queue; remove from pipeline                              |
| 2     | Unlink/release atee src pad; NULL a_queue; remove from pipeline                              |
| 3     | Unlink from comp_pad; release compositor pad; NULL/remove v_convert, v_decodebin             |
| 4     | Unlink from amix_pad; release audiomixer pad; NULL/remove a_resample, a_convert, a_decodebin |
| 5     | NULL webrtcbin; remove from pipeline                                                         |
| 6     | Free grid slot; remove from webrtcbins hash table                                            |
\end{verbatim}

\subsection{Signaling Protocol (Port IPC)}

Newline-delimited JSON over UNIX stdio. Stdin messages are processed synchronously
on the GLib main loop via `g\_io\_add\_watch`, serializing all pipeline mutations.

\subsubsection{Erlang to GStreamer (stdin)}

\begin{verbatim}
| Type                | Fields    | Effect                  |
|---------------------|-----------|-------------------------|
| `peer_joined`       | `peer_id` | Triggers `setup_peer()` |
| `sdp_answer`        | `peer_id`, `sdp` | Sets remote description on webrtcbin |
| `ice_candidate`     | `peer_id`, `candidate.sdpMLineIndex`, `candidate.candidate` | Adds ICE candidate |
| `peer_left`         | `peer_id` | Triggers six-stage `cleanup_peer()` |
| `exit`              | — | Sends EOS to mux; finalizes recording |
\end{verbatim}

\subsubsection{GStreamer to Erlang (stdout)}

\begin{verbatim}
| Type                | Fields    | Effect                  |
|---------------------|-----------|-------------------------|
| `sdp_offer`         | `peer_id`, `sdp` | Forwarded by broker to browser via Syn |
| `ice_candidate`     | `peer_id`, `candidate` | Forwarded to browser |
| `recording_started` | — | Pipeline reached PLAYING state |
\end{verbatim}

\subsection{Output Delivery Nuances}

\subsubsection{PTS/DTS Integrity}

HLS tees after `h264parse`, before `rtph264pay`. RTP payload/depayload cycles
corrupt timestamps; `mpegtsmux` is intolerant of discontinuities and MSE decoders
freeze permanently after the first affected segment.

\subsubsection{Disk-I/O Isolation}

30-second leaky queues (`max-size-time=30000000000`) on storage branches absorb
filesystem stalls, guaranteeing HLS/MP4 continuity independent of disk I/O jitter.

\subsubsection{Fragmented MP4}

`mp4mux fragment-duration=1000 streamable=true` interleaves `moof`/`mdat` atoms,
enabling progressive download and crash-resilient playback without a terminal `moov`.

\subsubsection{Dual-Encode}

`raw\_vtee` drives separate `x264enc` (WebRTC) and `x265enc` (HLS) encoder chains,
providing H.265 fidelity for capable clients without breaking WebRTC compatibility.

\subsection{Signal Handling and Graceful Shutdown}

\begin{verbatim}
GstElement *mux = gst_bin_get_by_name(GST_BIN(state.pipeline), "mux");
if (mux) {
    gst_element_send_event(mux, gst_event_new_eos());
    gst_object_unref(mux);
}
\end{verbatim}

The muxer finalizes index structures before the loop quits on bus EOS.

\subsection{Build and Evaluation}

\subsubsection{Dependencies (macOS)}

\begin{verbatim}
brew install gstreamer libnice libnice-gstreamer json-glib
\end{verbatim}

\subsubsection{Dependencies (Ubuntu / WSL2)}

\begin{verbatim}
sudo apt-get update
sudo apt-get install -y libgstreamer1.0-dev libgstreamer-plugins-base1.0-dev \
     libgstreamer-plugins-bad1.0-dev libjson-glib-dev pkg-config
\end{verbatim}

In PowerShell allow UDP traffic:

\begin{verbatim}
New-NetFirewallRule -DisplayName "RTP WebRTC Media UDP (WSL2)"
                    -Direction Inbound -Action Allow -Protocol UDP -LocalPort 49152-65535
\end{verbatim}

In `gst.c`:

\begin{verbatim}
g_object_set(webrtc, "latency", 50, "stun-server", "stun://stun.l.google.com:19302", NULL);
\end{verbatim}

\subsubsection{Compile}

\begin{verbatim}
cc -O3 c_src/gst.c -o priv/gst \
  $(pkg-config --cflags --libs \
    gstreamer-1.0 gstreamer-webrtc-1.0 gstreamer-sdp-1.0 json-glib-1.0)
\end{verbatim}

\subsubsection{Standalone Test}

\begin{verbatim}
mkdir -p /tmp/rtp-test && ./priv/gst /tmp/rtp-test ts
\end{verbatim}

\subsubsection{Launch via Erlang Node}

\begin{verbatim}
$ iex -S mix
$ open http://localhost:8081/app/login.htm
\end{verbatim}

\subsubsection{Mock Devices}

**Chrome**: `--use-fake-device-for-media-stream --use-fake-ui-for-media-stream`
**Safari**: Develop -> WebRTC -> Use Mock Capture Devices

\subsection{Architectural Invariants}

\begin{verbatim}
| Invariant | Mechanism |
|---|---|
| Continuous pipeline flow | `videotestsrc` + `audiotestsrc` on compositor/audiomixer `sink_0` |
| Thread-safe pipeline mutation | All mutations on GLib main loop via IO channel watch |
| Back-pressure isolation | Per-peer leaky queues (`leaky=2`, `max-size-time=1s`) |
| Z-order correctness | Background `zorder=1`; peers `zorder = grid_idx + 10` |
| PTS/DTS integrity | Tees after `h264parse`, before `rtph264pay` |
| Crash-resilient MP4 | `mp4mux streamable=true fragment-duration=1000` |
| Disk-I/O isolation | 30-second leaky queues on HLS/MP4 storage branches |
| Signal safety | `g_unix_signal_add()` dispatches to GLib main loop |
| Resource reclamation | Six-stage ordered cleanup on every peer departure |
| Grid capacity | Up to 16 simultaneous participants (4x4 spatial grid) |
\end{verbatim}

\subsection{Low Latency Config for Hardware}

Current Stable Configuration:

\begin{verbatim}
| Component            | Recommended Value | Notes                |
|----------------------|-------------------|----------------------|
| webrtcbin latency    | 100 ms            | Extectations         |
| Video jitter queue   | 200 ms            | After decode         |
| Audio jitter queue   | 250 ms            | After decode         |
| Outgoing v_queue     | 400 ms            | Critical             |
| Outgoing a_queue     | 450 ms            | Audio more sensitive |
| Pre-encoder queue    | 300 ms            | Before x264enc       |
| key-int-max          | 60                | Keep higher on Pi    |
| x264enc speed-preset | ultrafast         | Do not use superfast |
\end{verbatim}

\section{Conclusion}

\end{document}